\documentclass[landscape,twocolumn,leqno,a4paper,10pt]{report} %leqno affiche les numéros d'équations à gauche

\usepackage[a4paper,top=0.4cm,bottom=1cm,left=0.4cm,right=0.4cm,marginparwidth=0.5cm,footskip=0.6cm]{geometry}



%----------------- Codage de la langue et de la police--------------------------------------------------------------------------
\usepackage[french]{babel} 		% Permet de spécifier au compilateur que l'on désire taper son document en français.
\usepackage[utf8x]{inputenc} 	% Permet de taper directement à l'écran, dans le fichier .tex, les caractères accentués.
\usepackage[T1]{fontenc} 		% Permet juste d'utiliser une nouvelle norme LATEX concernant le codage des caractères.
%\usepackage{lmodern}		% Pack de police --> Pas besoin puisque j'utilise Fourier

%------------------------ Packages utiles-------------------------------------------------------------------------------------------

\usepackage{graphicx}	% Pour insérer des images
\usepackage{caption}		% Légende

\newcommand\hmmax{0}		% Alors c'est une astuce assez mystérieuse pour dépasser la limite de polices latex.
							% En effet, le package bm est lourd et avec tout ce que l'on met, on obtient un message d'erreur : Too many math alphabets used in version normal.			
\newcommand\bmmax{0}		% Cette astuce règle le problème, il faut tout de même mettre ces compteurs avant le package bm

\usepackage{amsmath,amssymb}	% Permettent d'obtenir des caractères mathématiques (de toute beauté !).
\usepackage{amsthm}		% Pour pouvoir faire les newthm
\usepackage{mathtools}	% Corrige des bugs dans amsmath et le complète
\usepackage{numprint}	% Pour les espacements dans les nombres : \numprint{10000} = 10 000
\usepackage{fourier}		% REMARQUE : [upright] impose aux majuscules d'être droite dans le mode mathématique 
\usepackage{pifont} 		% Importe une large gamme de symboles, on utilise \ding{numéro}
\usepackage{xargs} 		% Permet de définir des macros à plusieurs paramètres optionnels.
\usepackage{fancyhdr} 	% Package pour les entêtes et bas de page
\pagestyle{fancy} 			% Accompagne \usepackage{fancyhdr}
\usepackage{bm} 		% Permet de mettre en gras du contenu mathématiques via \bm{..}
\usepackage{mathrsfs} 	% Pour les lettres cursives
\usepackage{enumitem} 	% Pour pouvoir personnaliser les listes

\frenchbsetup{StandardLists=true} % Afin d'éviter tout conflit entre les extensions enumitem et french babel. 

\usepackage{sectsty}	% Pour colorer les titres
\usepackage{titlesec} 	% Pour personnaliser les titres
\usepackage{setspace} 	% Pour changer la taille des interlignes
       	               
\usepackage{array,multirow,makecell}
\usepackage{hlist}		% Package Christian Tellechea sur les listes numérotées gauche vers droite
\setdefaulthlist{left margin = 2mm, label sep= 0.2mm}	% réglages package hlist
				

%------------ Les packages de dessin sous pstricks : il faut changer la compilation de latex pour compiler ces dessins------------------
% Options\Configurer Texmaker\Compil rapide\ Latex+dvips+ps2pdf+Voir PDF

\usepackage[pdf]{pstricks} 	% Pour faire des dessins
\usepackage{pstricks-add,pst-plot,pst-tree,pst-node,pst-math,pst-eucl,pst-func,pst-3dplot}% ,pst-text,pst-eps,pst-fill,pst-grad
\usepackage{multido}		% Pour faire des multido
\usepackage{lastpage} 	% Pour le compteur de la dernière page, on appelle alors celle-ci avec lastage



\usepackage[colorlinks=true, 
			%allcolors=blue
			]{hyperref}		% ATTENTION : doit être importé après la majorité des packages sauf tabularx et breakrules (par exemple)
\hypersetup{			% Paramétrage des hyperliens
	colorlinks = true,   		% Colorise les liens
	breaklinks = true,		% Permet les retours à la ligne pour les liens trop longs
	urlcolor = blue,		% Couleur des hyperliens
	linkcolor = black,		% Couleur des liens internes aux documents (index, figures, tableaux, équations, ...)
	citecolor = blue		% Couleur des liens vers les références bibliographiques
	}




%%%%%%%%%%%%%%%%%%%%%%%%%%%%%%%%%%%%%%%%%%%%%%%%%%
%------------------------- Définition des commandes -----------------------------------------------------------------------------
%%%%%%%%%%%%%%%%%%%%%%%%%%%%%%%%%%%%%%%%%%%%%%%%%%


%--------------------- Raccourcis fontes--------------------------------------------------------------
% REMARQUE : le package fourier met la même police pour mathcal et mathscr, on remet donc mathcal normalement et on utilise la définition de fourier (un peu plus grasse) pour mathscr
\DeclareMathAlphabet{\mathcal}{OMS}{cmsy}{m}{n}	
\DeclareMathAlphabet{\mathscr}{FMS}{futm}{m}{n}

%---------------------------- Raccourcis des espaces--------------------------------------------------------------------------------------------
\newcommand\n{\noindent}
\newcommand\q{\quad}
\newcommand\qq{\quad\quad}

%--------------------- Raccourcis en mode mathématiques-------------------------------------------------------------------------------------
\newcommand{\et}{\text{ et }}
\newcommand{\ou}{\text{ ou }}
\newcommand{\SSI}{\Longleftrightarrow}
\let\si\relax
\newcommand{\si}{\text{ si }}
\newcommand{\sinon}{\text{ sinon }}
\renewcommand{\epsilon}{\varepsilon}		% epsilon beaucoup plus joli
\renewcommand{\geq}{\geqslant} 
\renewcommand{\leq}{\leqslant}
\newcommand{\e}{\mathrm{e}}
\renewcommand{\textbf}[1]{\begingroup\bfseries\mathversion{bold}#1\endgroup}		% Très important, permet de passer en gras même avec le mode mathématiques

\newcommand{\Mod}[1]{\allowbreak\if@display\mkern5mu\else\mkern5mu\fi\left[#1\right]}	% Pour modulo entre crochets [.] avec un espacement convenable 


%---------------------------------- Raccourcis calcul diff----------------------------------------------------------------------------------------
\renewcommand{\d}{\;\mathrm{d}}	% Pour le d droit de l'intégrale car ce n'est pas une variable...


%-------------------------------- Raccourcis des ensembles----------------------------------------------------------------------------
\newcommand\R{\mathbb{R}}
\newcommand\C{\mathbb{C}}
\newcommand\K{\mathbb{K}}
\newcommand\Q{\mathbb{Q}}
\newcommand\D{\mathbb{D}}
\newcommand\Z{\mathbb{Z}}
\newcommand\N{\mathbb{N}}
\newcommand\U{\mathbb{U}}


%------------------------------- Dimensions intelignes, items...-----------------------------------------------------------------------------------------------
\onehalfspacing	% Pour que l'interligne soit de 1.5. Pour que l'interligne soit de la valeur que l'on souhaite, de base : \setstretch{1}

\setlist[1]{font=\bfseries,itemsep=0.2mm,leftmargin=6mm}
\setlist[2]{font=\bfseries,itemsep=0.2mm,leftmargin=4mm}
\setenumerate[1]{font=\bfseries,itemsep=0.2mm,leftmargin=6mm}
\setenumerate[2]{font=\bfseries, label=\alph*., leftmargin=4mm}
%%\setenumerate[profondeur]{font=\bfseries, label=\ding{221} \color{blue}} permet de changer en a) b) les sous-items d'une liste enumerate
%% itemsep = correspond à l'interligne entre chaque item
% de base : labelsep=2mm et permet de changer la distance entre l’étiquette et le texte des éléments
% leftmargin =4mm spécifie la marge à gauche de la liste
\AtBeginDocument{\renewcommand{\labelitemi}{\textbullet}}	% Pour que les puces par défaut soit des points.
%\AtBeginDocument{\renewcommand{\labelitemi}{\ding{86}}}	% Pour que les puces par défaut soit des étoiles.
\AtBeginDocument{\renewcommand{\labelitemii}{$\ast$}}		% Pour que les puces par défaut pour une sous-liste soit des astérisques.
\newcommand{\itemmm}{\item[$\ast$]}				% \itemmm devient l'item \ast
\newcommand{\itemmmm}{\item[\footnotesize$\circledcirc$\normalsize]}
\newcommand{\itemmmmm}{\item[\footnotesize$\astrosun$\normalsize]}
\newcommand{\itemmmmmm}{\item[\footnotesize$\multimap$\normalsize]}
\newcommand{\itemmmmmmmmm}{\item[$\divideontimes$]}
\newcommand{\etoile}{\item[$\star$]}
\newcommand{\etoilee}{\item[$\bigstar$]}
\newcommand{\losange}{\item[$\diamondsuit$]}
\newcommand{\carre}{\item[\LARGE$\centerdot$]}
\newcommand{\carreombre}{\item[\ding{112}]}

\setlength{\columnsep}{30pt} 		% Permet d'écarter les deux colonnes du twocolumn de 30pt, pratique pour écarter le texte de la ligne centrale
\setlength{\columnseprule}{0.5pt}	% Affiche une ligne verticale entre chaque item et entre chaque page du twocolumn
						% \columnseprule=0cm 	% Permet d'enlever le trait vertical entre chaque colonne
%\renewcommand\tabenumsep{0.5cm\hskiplem}	% Ecarte de 0.5cm les items de tabenum, sinon ils sont trop collés...


\newlength{\tagmarginsep} 		% On crée une longueur
\setlength{\tagmarginsep}{0cm}		% Permet de régler la position d'un tag pour une équation mettre avant l'équation : \setlength{\tagmarginsep}{x cm} pour que le tag se rapproche de l'équation
\everydisplay{\displayindent=\tagmarginsep \displaywidth=\dimexpr\linewidth-2\tagmarginsep}	% Va avec ce qu'il y a juste avant


%-------------------- Définition des couleurs-----------------------------------------------------------------------------------------

\definecolor{violetfonce}{RGB}{126,62,142}
\definecolor{magentafonce}{RGB}{128, 0, 128}
\newrgbcolor{magentafonce}{.5 0 .5}			% Une autre façon de définir une couleur, permet surtout d'utiliser \magenta
\definecolor{violet1}{cmyk}{0.5,1,0,0}
\definecolor{violet2}{cmyk}{0.24998, 1, 0, 0}
\definecolor{violet3}{cmyk}{0, 1, 0.25, 0}



\newcommand\tred{\textcolor{red}}
\newcommand\tredf{\textcolor{rougefonce}}
\newcommand\tdred{\textcolor{deepred}}
\newcommand\tblue{\textcolor{blue}}
\newcommand\tbluef{\textcolor{bleufonce}}
\newcommand\tgreen{\textcolor{vertfonce}}
\newcommand\torange{\textcolor{orange}}
\newcommand\tmagenta{\textcolor{magentafonce}}


%------------------ Changement paramètres des parties, sous parties---------------------------------------------------------


\sectionfont{\color{violet1}}			% Utilisation du package sectsty
\subsectionfont{\color{violet2}}
\subsubsectionfont{\color{violet3}}


%--------------------------------- Raccourcis probabilités--------------------------------------------------------------------------------
\newcommand{\Prob}{\mathbb{P}}
\newcommand{\E}{\mathbb{E}}
\newcommand{\normale}[2]{\mathcal{N}\left(#1;#2 \right)}	% Pour la loi Normale : 	$\normale{\mu}{\sigma^2}$
\newcommand{\Unif}[1]{\mathcal{U}\left(#1\right)}			% Pour la loi uniforme :	$\Unif{[a;b]}$
\newcommand{\Ber}[1]{\mathcal{B}\left(#1\right)}			% Pour la loi de Bernoulli : $\Ber{p}$
\newcommand{\Bino}[2]{\mathcal{B}\left(#1;#2\right)}		% Pour la loi Binomiale :	$\Bino{n}{p}$
\newcommand*{\binomial}[2]{% 
  \ensuremath{ 
    \begin{pmatrix} 
      #2\\ 
      #1
    \end{pmatrix}}}				% Pour le coefficient binomial : $\binom{k}{n}$
\renewcommand{\binom}{\binomial} 	% La commande \binom existait déjà mais n'était pas pratique



%----------------------------------- Opérateurs---------------------------------------------------------------------------------------------------

\newcommand{\barre}[1]{\overline{#1}}  									% Barre au dessus du texte, pratique pour le conjugué des complexes
\newcommandx{\Sum}[3][3=\white]{{\displaystyle\sum\limits_{\substack{#1 #3}}^{#2}}}		% $\Sum{k=0}{n}[<option>]$ où [<option>] est placée à la suite de k=0 si l'on tape de façon classique, et en-dessous si l'on met [\\ k\neq k_0]
															% On peut juste taper : $\Sum{k=0}{n}$ et cela fonctionne
															% \displaystyle impose de gérer les tailles comme en mode mathématique : ici on aura donc toujours de grande somme que ce soit en math comme dans le texte
\newcommandx{\Prod}[3][3=\white]{{\displaystyle\prod\limits_{\substack{#1 #3}}^{#2}}}		% $\Prod{k=0}{n}[<option>]$	idem pour l'option
\newcommand{\Int}[4]{{\displaystyle\int_{#1}^{#2} #3  \d #4}}   					% $\Int{0}{\pi}{\cos(t)}{t}$
\newcommandx{\Lim}[3][3=\white]{\lim\limits_{\substack{#1 \to #2 #3}}}		% $\Lim{x}{0}[<option>]$	idem pour l'option


%------------------------------------ Délimiteurs--------------------------------------------------------------------------------------------------
%\newcommand{\abs}[1]{\left\lvert #1 \right\rvert}			% Valeur absolue | |
%\newcommand{\norm}[1]{\left\lVert #1 \right\rVert}			% Doubles barres verticales : norme d'un vecteur || ||

%--------------------------------------- Vecteurs----------------------------------------------------------------------------------------------------
 
\newcommand{\Oij}{\left(O;\vec{\imath},\vec{\jmath}\right)}
\newcommand{\Oijk}{\left(O;\vec{\imath},\vec{\jmath},\vec{k}\right)}
\newcommand{\Ouv}{\left(O;\vec{u},\vec{v}\right)}
\newcommand{\Ouvw}{\left(O;\vec{u},\vec{v},\vec{w}\right)}

%------------------------------------ Systèmes d'équations-------------------------------------------------------------------------------------------
% Remarque \syst et \syste ne font pas d'alignement mais ils sont bien, si l'on veut de l'alignement, on utilisera \systeme du package "systeme"
\newcommand*{\syst}[2]{%			% Systèmes 2 éq. à ? inconnues : 	\syst{ax+by+c=0}{a'x+b'y+c'=0}
  \left\{\begin{array}{ll} 
        #1\\ 
        #2 
      \end{array}  
    \right. }
 
\newcommand*{\syste}[3]{% 			% Systèmes 3 éq. à ? inconnues : 	\syste{ax+by+c=0}{a'x+b'y+c'=0}{a''x+b''y+c''=0}
  \left\{\begin{array}{ll} 
        #1\\ 
        #2\\
        #3 
      \end{array}  
    \right.}



%------------------------------------------ Autre----------------------------------------------------------------------------------------------------------
\newcommand{\boite}[2]{\noindent\parbox[]{#1 cm}{#2} }  									% Boite autre façon de faire minipage en gros


%---------------------------- En-têtes et bas de pages----------------------------------------------------------------------------
% A chaque fois : *head* pour le haut et *foot* pour le bas
% L = left, C = center, R = right 
% L = toutes les pages, LE = toutes les pages paires, LO = toutes les pages impaires, idem CE, CO, ...


\newcommand{\fcours}[1]{								% f pour fancy				   
	\pagestyle{fancy}									% Efface le style qui a pu être mis avant
	\renewcommand{\headrulewidth}{0pt} 
	\renewcommand{\footrulewidth}{0.5pt} 
	%\fancyhead[L]{}									
	%\fancyhead[C]{}
	%\fancyhead[R]{}									
	\fancyfoot[L]{\vspace{-1mm}\footnotesize \textit{#1}} 			% Classe : <Seconde> 
	\fancyfoot[C]{\vspace{-1mm}\footnotesize \thepage/\pageref{LastPage}}
	\fancyfoot[R]{\vspace{-1.8mm}\textit{\footnotesize C. Marguerite}}	% N'apparaîtra pas si on met les marges avec moins de bottom=1.3cm dans le package geometry
 }												% Fin de la commande
					%% Mettre au début du document : \fcours{Seconde} par exemple
					%% \usepackage[a4paper,top=0.4cm,bottom=1cm,left=0.4cm,right=0.4cm,marginparwidth=0.5cm,footskip=0.6cm]{geometry} 




%---------------------------- Réglages chapitres + table des matières------------------------------------------------------
\titleformat{\part}[display]
{\centering\normalfont\Huge\bfseries}
{\partname}
{20pt}
{\rule{6cm}{1pt}\raisebox{-2.1pt}[10pt][10pt]{\quad\decofourleft\decoone\decofourright\quad}\rule{6cm}{1pt} ~\\\vspace{20pt}}

\makeatletter
\renewcommand\chapter{\global\@topnum\z@
                    \@afterindentfalse
                    \secdef\@chapter\@schapter}
\def\@topnewpage[#1]{#1}
\makeatother										% Pour avoir le chapitre sur une demie page

\titleformat{\chapter}{\centering\normalfont\Large\bfseries\scshape\color{magentafonce}}{}{1em}{}
\titlespacing*{\chapter}{0pt}{-18pt plus 0pt minus 2pt}{0pt plus 2pt minus 2pt}		% Réglages espacements (à gauche, au-dessus, en-dessous) de la commande \chapter

\makeatletter
\newcommand{\twocolumntoc}{		
  \section*{\huge\centering\contentsname
    \@mkboth{%
      \MakeUppercase\contentsname}{\MakeUppercase\contentsname}}%
  \@starttoc{toc}%
}
\makeatother										% Permet de mettre la table des matières en twocolumn 
												% Mettre après \begin{doc..} \twocolumntoc \clearpage 		% Remplace \tableofcontents par \twocolumntoc en twocolumn

\makeatletter
\renewcommand{\thesection}{\@arabic\c@section}
\makeatother										% Enlève la numérotation des chapitres devant les sections

\setcounter{tocdepth}{3} 								% Permet d'afficher dans la table des matières suivant une certaine profondeur : -1=part, 0=chapter, 1=section, ...
\setcounter{secnumdepth}{3}							% Pour avoir la numérotation des subsubsections avec les classes book et report

\newcommand{\nocontentsline}[3]{}
\let\oldaddcontentsline\addcontentsline
\newcommand{\tocless}[2]{\let\addcontentsline\nocontentsline #1{#2}\let\addcontentsline\oldaddcontentsline}	% Pour retirer une section de la table des matières : \tocless\section{...}





\setstretch{1}
\begin{document}
\fcours{Concours général}


\chapter{Concours général des lycées - Session 2016}

\section{Problème 1 : Sommes de cubes}

\n Si $n$ est un entier, on appelle cube de $n$ l'entier $n^{3}$.\\
Dans tout le problème, on note :
\begin{itemize}
	\itemmm $S$ l'ensemble des entiers strictement positifs qui peuvent se décomposer en une somme de cubes d'entiers strictement positifs deux à deux distincts;
	\itemmm $S_{0}$ l'ensemble des entiers strictement positifs qui peuvent se décomposer en une somme de cubes d'entiers pairs strictement positifs deux à deux distincts;
	\itemmm $S_{1}$ l'ensemble des entiers strictement positifs qui peuvent se décomposer en une somme de cubes d'entiers impairs strictement positifs deux à deux distincts.
\end{itemize}

\n Par exemple, 8 et 190 sont dans $S$ car $8=2^{3}$ et $190=1^{3}+4^{3}+5^{3}$; $216$ et 1072 sont dans $S_{0}$ car $216=6^{3}$ et $1072=2^{3}+4^{3}+10^{3}$; $125$ et 2568 sont dans $S_{1}$ car $125=5^{3}$ et $2568=1^{3}+3^{3}+7^{3}+13^{3}$.\\
L'objectif du problème est de démontrer que tout entier suffisamment grand appartient à $S$.

\begin{enumerate}
	\item Montrer que 2016 appartient à $S_{0}$.
	\item \begin{enumerate}
		\item Montrer que, pour tout réel $x \geq 5$, on a $(2 x+1)^{3} \leq 2(2 x-1)^{3}$.
		\item Soit $k$ un entier supérieur ou égal à $5$. Montrer, pour tout entier $p \geq k$,
		$$(2 p+1)^{3} \leq(2 k-1)^{3}+\sum_{j=k}^{p}(2 j-1)^{3}$$
		où on rappelle que $\Sum{j=k}{p}(2 j-1)^{3}$ désigne la somme $(2 k-1)^{3}+(2 k+1)^{3}+\cdots+(2 p-1)^{3}$.
	\end{enumerate}
\end{enumerate}
On rappelle que si $t, u$ et $v$ sont des entiers, la notation $t \equiv u(v)$ signifie que $v$ divise $u-t$.

\begin{enumerate}\setcounter{enumi}{2}
	\item Montrer qu'il existe 288 entiers $s_{1}, \ldots, s_{288}$ appartenant à $S_{1}$ tels que, pour tout $i$ dans $\{1, \ldots, 288\}$, $s_{i} \equiv i \Mod{288}$.
\end{enumerate}
Dans la suite du problème, on fixe des entiers $s_{1}, \ldots, s_{288}$ appartenant à $S_{1}$ tels que $s_{i} \equiv i\Mod{288}$ pour tout $i$, et on note $m$ le plus grand des nombres $s_{i}$ :
$$m=\max \left(s_{1}, s_{2}, \ldots, s_{288}\right)$$
On rappelle que $n$ réels $u_{1}, \ldots, u_{n}$ sont dits en progression arithmétique de raison $r$ si $u_{i+1}-u_{i}=r$ pour tout entier $i$ tel que $1 \leq i<n$.
\begin{enumerate}\setcounter{enumi}{3}
	\item Soit $n$ un entier tel que $288 n \geq m$, et soit $u_{1}, \ldots, u_{n}$ des entiers naturels en progression arithmétique de raison 288 . Montrer que tout entier de l'intervalle $\left[m+u_{1}, 288 n+u_{1}\right]$ peut s'écrire sous la forme $s_{i}+u_{j}$, avec $1 \leq i \leq 288$ et $1 \leq j \leq n$.
	\item On admet la relation, pour tout réel $x$,
	$$(2 x+12)^{3}+(2 x+4)^{3}+(2 x+2)^{3}-(2 x+10)^{3}-(2 x+8)^{3}-(2 x)^{3}=288$$
	\begin{enumerate}
		\item Montrer qu'il existe un entier $u$ tel que $u$, $u+288$ et $u+576$ appartiennent à $S_{0}$.
		\item Montrer que, pour tout entier $n \geq 2$, il existe $n$ éléments dans $S_{0}$ en progression arithmétique de raison 288.
	\end{enumerate}
	\item Soit $k$ un entier supérieur ou égal à 5 tel que $(2 k+1)^{3}>m$.
	\begin{enumerate}
		\item Montrer qu'il existe un entier $N \geq 1$ tel que tout entier de l'intervalle $\left[N, N+2(2 k-1)^{3}\right]$ puisse s'écrire sous la forme $s_{i}+u$, avec $1 \leq i \leq 288$ et $u \in S_{0}$.
		\item Montrer que tout entier supérieur ou égal à $N$ appartient à $S$.\\
		\textit{Pour tout entier $p \geq k$, on pourra examiner le cas des entiers de l'intervalle $\left[N, N_{p}\right]$, où}
		$$N_{p}=N+(2 k-1)^{3}+\sum_{j=k}^{p}(2 j-1)^{3}$$
	\end{enumerate}
\end{enumerate}

\newpage

\section{Problème 2 : La rangée d'arbres qui cache la forêt}

\n Un observateur se trouve dans une forêt parfaitement régulière, dans laquelle les troncs d'arbres ont tous même diamètre. On ne tient pas compte de la hauteur, ce qui permet de se ramener à un problème dans le plan, qu'on rapporte à un repère orthonormé. L'observateur se trouve à l'origine du repère et les troncs d'arbres sont représentés par des cercles de même rayon $R>0$, centrés aux points de coordonnées $(a, b)$, où $a$ et $b$ sont des entiers impairs. Il s'agit d'étudier ce que voit l'observateur.\\
Pour des raisons de symétrie, on peut se restreindre au quart de plan $x>0$, $y>0$. On dira que l'observateur << voit à travers la forêt >> s'il existe une demi-droite issue de l'origine, contenue dans le quart de plan considéré et ne rencontrant aucun des cercles.

\begin{center}
\psset{xunit=1cm,yunit=1cm,arrowsize=4pt,plotpoints=1000,algebraic=true}
\begin{pspicture*}(-0.5,-0.5)(8,8)
\psaxes[Dx = 1,Dy =1,labelsep=0.8mm,linewidth=.75pt,ticksize=-2pt 2pt,showorigin=false]{->}(0,0)(-0.5,-0.5)(8,8)
\uput[dl](0,0){$0$}
\def\F{3.5*x}
\def\G{2*x}
\psplot{0}{3}{\F}
\psplot{0}{4}{\G}
\multido{\i=1+2}{4}{\multido{\n=1+2}{4}{\pscircle(\i,\n){0.3}}}
\uput[dr](2,7){$\mathscr{D}$}
\uput[dr](3.5,7){$\mathscr{D}'$}
\psdot(0.79,2.78)
\end{pspicture*}
\end{center}


\n Par exemple, dans la figure ci-dessus, la demi-droite $\mathscr{D}$ rencontre un arbre, mais pas la demi-droite $\mathscr{D}^{\prime}$. Pour $m \in] 0,+\infty[$, on note $\mathscr{D}_{m}$ la demi-droite définie par les conditions $y=m x$ et $x>0$.\\

\n \textit{Dans ce problème, on admettra que si $m$ est un nombre irrationnel positif, alors, pour tout réel strictement positif $\varepsilon$, il existe deux entiers naturels impairs $a$ et $b$ tels que $|b-m a| \leq \varepsilon$}.

\begin{enumerate}
	\item Soient $a$, $b$, $m$ des réels strictement positifs. Montrer que $\mathscr{D}_{m}$ rencontre le cercle de rayon $R>0$ centré en $(a, b)$ si et seulement si $|b-m a| \leq R \sqrt{1+m^{2}}$.
	\item En déduire que si $m$ est irrationnel, alors $\mathscr{D}_{m}$ rencontre un arbre.
	\item On suppose que $m=\frac{b}{a}$, où $a$ et $b$ sont des entiers naturels non nuls premiers entre eux.
	\begin{enumerate}
		\item On suppose que $a$ et $b$ sont impairs. La demi-droite $\mathscr{D}_{m}$ rencontre-t-elle un arbre ?
		\item On suppose que $a$ et $b$ sont de parités différentes et que $\mathscr{D}_{m}$ rencontre un arbre.\\
		Montrer que $1 \leq R \sqrt{a^{2}+b^{2}}$.
	\end{enumerate}
	\item En déduire que si toutes les demi-droites $\mathscr{D}_{m}$, avec $m>0$, rencontrent un arbre, alors $R \geq \frac{1}{\sqrt{5}}$.
	\item On suppose réciproquement que $R \geq \frac{1}{\sqrt{5}}$. Montrer que toute demi-droite $\mathscr{D}_{m}$, avec $m>0$, rencontre un arbre planté en $(\alpha, 1)$ ou en $(1, \alpha)$, où $\alpha$ est un entier naturel impair.
\end{enumerate}

\n On appelle première rangée d'arbres l'ensemble des arbres plantés aux points $(\alpha, 1)$ ou $(1, \alpha)$, où $\alpha$ est un entier naturel impair.

\begin{enumerate}\setcounter{enumi}{5}
	\item Conclure que si l'observateur voit à travers la première rangée, alors il voit à travers la forêt.
\end{enumerate}

\newpage


\section{Problème 3 : Allons dans $\mathbb{C}$}


\n Dans tout le problème, $j$ désigne le nombre complexe $\e^{2 \pi i / 3}$.\\
La probabilité d'un évènement $A$ est notée $\Prob(A)$.

\begin{enumerate}
	\item \begin{enumerate}
		\item Vérifier que $j^{3}=1$ et $1+j+j^{2}=0$.
		\item Que peut-on dire du triangle dont les sommets ont pour affixes $1$, $j$, $j^{2}$ ?
		\item Montrer que si $a$, $b$, $c$ sont des nombres réels, alors $a+b j+c j^{2}=0$ si et seulement si $a=b=c$.
	\end{enumerate}
\end{enumerate}

\n On lance un dé équilibré (six faces numérotées de 1 à 6 ). On note $F$ la variable aléatoire donnant le nombre obtenu, et on note $Z$ la variable aléatoire $j^{F}$.

\begin{enumerate}\setcounter{enumi}{1}
	\item Montrer que $Z$ est à valeurs dans $\left\{1, j, j^{2}\right\}$ et que $\Prob(Z=1)=\Prob(Z=j)=\Prob\left(Z=j^{2}\right)=\frac{1}{3}$.
\end{enumerate}

\n On considère un entier $n \geq 1$ et on lance le dé $n$ fois (lancers indépendants). On note $F_{k}$ le résultat du $k$-ième lancer et $Z_{k}=j^{F_{k}}$. On note $S_{n}=Z_{1}+\cdots+Z_{n}$ et $p_{n}=\Prob\left(S_{n}=0\right)$. On note $U_{n}$ la variable aléatoire qui donne le nombre d'entiers $k \in[1, n]$ tels que $Z_{k}=1$; on note $V_{n}$ celle qui donne le nombre d'entiers $k \in[1, n]$ tels que $Z_{k}=j$ et $W_{n}$ celle qui donne le nombre d'entiers $k \in[1, n]$ tels que $Z_{k}=j^{2}$.

\begin{enumerate}\setcounter{enumi}{2}
	\item \begin{enumerate}
		\item Déterminer $U_{n}+V_{n}+W_{n}$.
		\item Montrer que $S_{n}=U_{n}+j V_{n}+j^{2} W_{n}$.
		\item Montrer que $S_{n}=0$ si et seulement si $U_{n}=V_{n}=W_{n}$.
		\item En déduire que si $n$ n'est pas multiple de 3, alors $p_{n}=0$.
	\end{enumerate}
	\item On suppose qu'il existe un entier naturel non nul $m$ tel que $n=3m$.
	\begin{enumerate}
		\item Montrer que la variable aléatoire $U_{n}$ suit une loi binomiale dont on précisera les paramètres.
		\item En déduire que $\Prob\left(U_{n}=m\right)=\binom{m}{3m} \frac{2^{2 m}}{3^{3 m}}$.
	\end{enumerate}
	On note $\Prob_{U_{n}=m}\left(V_{n}=m\right)$ la probabilité conditionnelle de $V_{n}=m$ sachant $U_{n}=m$.
	\begin{enumerate}\setcounter{enumii}{2}
		\item Montrer $\Prob_{U_{n}=m}\left(V_{n}=m\right)=2^{-2 m}\binom{m}{2m}$.
		\item En déduire $p_{3 m}=3^{-3 m}\binom{m}{3m}\binom{m}{2m}$.
	\end{enumerate}
\end{enumerate}

\n La question précédente, combinée à une expression classique des coefficients binomiaux, entraine pour $m$ entier naturel non nul la relation suivante, qu'on ne demande pas de démontrer :
$$\frac{p_{3 m+3}}{p_{3 m}}=\frac{(3 m+2)(3 m+1)}{9(m+1)^{2}}$$

\begin{enumerate}\setcounter{enumi}{4}
	\item Pour tout entier $m \geq 1$, montrer que $\frac{m}{m+1} \leq \frac{p_{3 m+3}}{p_{3 m}}$ et en déduire que $p_{3 m} \geq \frac{2}{9 m}$.
\end{enumerate}

\n Soit $X_{n}$ la variable aléatoire qui donne le nombre d'entiers $k \in[1, n]$ tels que $S_{k}=0$.

\begin{enumerate}\setcounter{enumi}{5}
	\item \begin{enumerate}
		\item Déterminer des variables de Bernoulli $Y_{k}$, avec $1 \leq k \leq n$, telles que $X_{n}=Y_{1}+\cdots+Y_{n}$.
		\item On note $\E\left(X_{n}\right), \E\left(Y_{1}\right), \ldots, \E\left(Y_{n}\right)$ les espérances de $X_{n}, Y_{1}, \ldots, Y_{n}$.\\
		En admettant que $\E\left(X_{n}\right)=\E\left(Y_{1}\right)+\cdots+\E\left(Y_{n}\right)$, montrer que $\E\left(X_{n}\right)=p_{1}+\cdots+p_{n}$.
		\item En déduire que $\Lim{n}{+\infty}\E\left(X_{n}\right)=+\infty$.
	\end{enumerate}
\end{enumerate}

\n Soit $q_{n}$ la probabilité que l'un des $S_{k}$ soit nul pour $1 \leq k \leq n$, c'est-à-dire $q_{n}=\Prob\left(X_{n}>0\right)$.\\
L'objectif de la question suivante est de montrer que la suite $\left(q_{n}\right)$ converge vers $1$.

\begin{enumerate}\setcounter{enumi}{6}
	\item \begin{enumerate}
		\item Montrer que la suite $\left(q_{n}\right)$ converge vers un réel $q$ et que $q_{n} \leq q \leq 1$ pour tout $n$.
		\item Pour $r$, $n$ entiers naturels non nuls, montrer que $\Prob\left(X_{n} \geq r\right) \leq q^{r}$.
		\item En déduire, pour tout entier $n \geq 1$, que $\E\left(X_{n}\right) \leq q+q^{2}+\cdots+q^{n}$.
		\item Conclure.
	\end{enumerate}
\end{enumerate}

\end{document}